\chapter{تعاریف اولیه}

فصل نخست مسئلۀ تحلیل احساسات ضمنی و نمای کلی سامانۀ پژوهش را معرفی کرد. برای تشریح دقیق ادبیات تحقیق و روش پیشنهادی، لازم است اصطلاحات اصلی با معنایی یکسان به کار روند. این فصل مفاهیم پایۀ تحلیل احساسات، مدل‌های زبانی، استدلال چندمرحله‌ای و ارزیابی طبقه‌بندی را تعریف می‌کند. تمرکز فصل بر تعریف‌ها و رابطۀ میان مفاهیم است؛ جزئیات پژوهش‌های پیشین در فصل سوم و نحوۀ پیاده‌سازی سامانه در فصل چهارم ارائه می‌شوند.

\section{تحلیل احساسات}

تحلیل احساسات یکی از مسائل پردازش زبان طبیعی است که نگرش یا ارزیابی بیان‌شده در متن را شناسایی می‌کند. خروجی این تحلیل می‌تواند قطبیت، شدت احساس یا رابطۀ نگرش با یک هدف مشخص باشد. در این پایان‌نامه، قطبیت نسبت به یک هدف معین اهمیت دارد؛ بنابراین، وجود یک واژۀ مثبت یا منفی در جمله برای تعیین پاسخ کافی نیست.

تحلیل احساسات با تشخیص هیجان یکسان نیست. هیجان به حالت‌هایی مانند شادی، خشم، ترس یا اندوه اشاره دارد. احساس در معنای مورد نظر این پژوهش، ارزیابی مطلوب، نامطلوب یا خنثی نسبت به یک موجودیت است. همچنین، ذهنی‌بودن یک جمله لزوماً قطبیت آن را تعیین نمی‌کند. ممکن است جمله‌ای ذهنی باشد، اما ارزیابی روشنی دربارۀ هدف مورد نظر نداشته باشد.

\subsection{سطوح تحلیل احساسات}

سطح تحلیل مشخص می‌کند که قطبیت برای کدام واحد زبانی یا معنایی پیش‌بینی می‌شود. سه سطح رایج عبارت‌اند از:

\begin{itemize}
    \item \textbf{سطح سند:} یک برچسب برای کل متن تعیین می‌شود. این سطح معمولاً فرض می‌کند که سند دربارۀ یک موضوع اصلی و دارای یک نگرش غالب است.
    \item \textbf{سطح جمله:} هر جمله جداگانه طبقه‌بندی می‌شود. وجود چند هدف یا چند نگرش متفاوت در یک جمله می‌تواند این سطح را با ابهام روبه‌رو کند.
    \item \textbf{سطح هدف یا جنبه:} قطبیت نسبت به هدف یا ویژگی مشخصی از آن تعیین می‌شود. در این سطح، یک جمله می‌تواند برای دو هدف، دو برچسب متفاوت داشته باشد.
\end{itemize}

سطح سوم برای این پژوهش مناسب‌تر است، زیرا ورودی شامل متن و هدف تعیین‌شده است. در یک صورت‌بندی عمومی، متن با $x$، هدف با $t$ و برچسب قطبیت با $y$ نمایش داده می‌شود. تابع طبقه‌بندی نگاشت زیر را می‌آموزد یا اجرا می‌کند:

\begin{equation}
    f(x,t) \longrightarrow y, \qquad y \in \mathcal{Y}.
    \label{eq:target-sentiment-function}
\end{equation}

مجموعۀ $\mathcal{Y}$ فضای برچسب‌ها است. هدف $t$ نیز دامنۀ ارزیابی را محدود می‌کند. در نتیجه، قطبیت موجودیت‌های دیگر نباید به هدف نسبت داده شود.

\subsection{قطبیت احساس}

قطبیت جهت کلی ارزیابی را نشان می‌دهد. فضای برچسب این پایان‌نامه از سه کلاس تشکیل شده است:

\begin{itemize}
    \item \textbf{مثبت:} متن، هدف را مطلوب، مفید یا سازگار با انتظار نشان می‌دهد.
    \item \textbf{منفی:} متن، هدف را نامطلوب، زیان‌آور یا ناسازگار با انتظار نشان می‌دهد.
    \item \textbf{خنثی:} شواهد کافی برای ارزیابی مثبت یا منفی هدف وجود ندارد.
\end{itemize}

برچسب خنثی فقط به معنای نبود واژۀ احساسی نیست. یک متن بدون واژۀ احساسی می‌تواند نگرش ضمنی داشته باشد. از سوی دیگر، وجود واژۀ احساسی نیز لزوماً قطبیت هدف را تعیین نمی‌کند. ممکن است آن واژه به موجودیتی دیگر اشاره کند یا در ساختاری مانند نفی و مقایسه به کار رفته باشد.

در این پژوهش، خنثی‌بودن با ابهام نیز تفاوت دارد. ابهام به نبود قطعیت در تفسیر اشاره می‌کند، اما خنثی یک برچسب معنایی است. سامانۀ طبقه‌بندی باید با وجود ابهام، یکی از برچسب‌های تعریف‌شده را انتخاب کند. نحوۀ مدیریت نمونه‌های مبهم به دستور مسئله و سیاست تصمیم‌گیری وابسته است.

\section{تحلیل احساسات مبتنی بر جنبه}

تحلیل احساسات مبتنی بر جنبه\LTRfootnote{Aspect-Based Sentiment Analysis (ABSA)} نگرش به موجودیت یا ویژگی مشخص را بررسی می‌کند و از تحلیل کل سند جزئی‌تر است. وظیفۀ چهارم \lr{SemEval 2014} معیار شناخته‌شدۀ این مسئله است \cite{pontiki2014semeval}.

\subsection{هدف و جنبه}

\textbf{هدف}\LTRfootnote{Target} موجودیتی است که قطبیت نسبت به آن سنجیده می‌شود. هدف می‌تواند نام محصول، خدمت، سازمان، ویژگی یا هر عبارت مشخص در متن باشد. \textbf{جنبه}\LTRfootnote{Aspect} ویژگی یا بُعدی از هدف است که ارزیابی به آن مربوط می‌شود. برای نمونه، «پشتیبانی» می‌تواند هدف باشد و «زمان پاسخ‌گویی» جنبۀ مورد ارزیابی آن باشد.

هدف و جنبه همیشه یکسان نیستند. هدف ممکن است به‌صورت صریح در متن ظاهر شود، اما جنبه از یک رویداد یا پیامد استنباط شود. همچنین، یک هدف می‌تواند چند جنبه داشته باشد. تعیین درست این رابطه مانع انتقال احساس از یک ویژگی به ویژگی دیگر می‌شود.

مفهوم \textbf{دستۀ جنبه}\LTRfootnote{Aspect Category} نیز با عبارت جنبه تفاوت دارد. دستۀ جنبه یک رده‌بندی کلی مانند «کیفیت خدمت» است، درحالی‌که عبارت جنبه ممکن است واژه‌ای مشخص در متن باشد. در تحلیل هدف‌محور این پایان‌نامه، هدف ورودی معلوم است، اما شناسایی جنبۀ مرتبط می‌تواند بخشی از فرایند استدلال باشد.

\subsection{نظر و قطبیت}

\textbf{نظر}\LTRfootnote{Opinion} محتوایی است که ارزیابی نویسنده را نسبت به هدف یا جنبه منتقل می‌کند. در حالت صریح، این محتوا معمولاً با یک عبارت ارزیابانه بیان می‌شود. در حالت ضمنی، رویداد، پیامد، توصیه یا مقایسه می‌تواند نقش نظر را داشته باشد.

برای نمونه، در جملۀ «پشتیبانی با نخستین تماس مشکل را حل کرد»، عبارت «مشکل را حل کرد» یک رویداد توصیفی است. با توجه به هدف «پشتیبانی» و انتظار معمول از خدمت، این رویداد می‌تواند نظر مثبت را منتقل کند. بنابراین، قطبیت فقط ویژگی یک واژه نیست؛ حاصل رابطۀ متن، هدف، جنبه و نظر است.

این رابطه را می‌توان با چهارگانۀ مفهومی $(t,a,o,y)$ نمایش داد. در این چهارگانه، $t$ هدف، $a$ جنبه، $o$ نظر و $y$ قطبیت است. همۀ مؤلفه‌ها لزوماً به صورت عبارت صریح در متن حضور ندارند. در تحلیل احساسات ضمنی، جنبه یا نظر ممکن است متغیری نهفته باشد که باید از بافت استنتاج شود.

\section{تحلیل احساسات ضمنی}

تحلیل احساسات ضمنی به تعیین نگرشی می‌پردازد که بدون نشانگر ارزیابانۀ مستقیم منتقل شده است. پژوهش‌های مرتبط، نیاز به دانش عمومی، درک بافت و استنتاج چندمرحله‌ای را از ویژگی‌های اصلی این مسئله می‌دانند \cite{fei2023thor,duan2024isa}. در این وضعیت، سامانه باید رابطۀ میان رویداد توصیف‌شده و مطلوبیت آن برای هدف را کشف کند.

\subsection{تفاوت احساس صریح و ضمنی}

در احساس صریح، عبارت‌هایی مانند صفت ارزیابانه، قید شدت یا فعل نگرشی مسیر نسبتاً مستقیمی به قطبیت ایجاد می‌کنند. در احساس ضمنی، قطبیت از پیامد یا معنای کاربردی یک گزاره به‌دست می‌آید. این تمایز به حضور یک واژۀ منفرد محدود نمی‌شود و باید نسبت به هدف سنجیده شود.

یک جمله می‌تواند هم‌زمان احساس صریح و ضمنی داشته باشد. ممکن است نگرش به یک هدف صریح باشد، اما نگرش به هدفی دیگر از رابطه‌ای غیرمستقیم استنباط شود. همچنین، نفی، کنایه یا نقل قول می‌توانند ظاهر واژگانی را با معنای نهایی ناسازگار کنند. ازاین‌رو، صریح یا ضمنی‌بودن یک ویژگی وابسته به متن، هدف و بافت است.

احساس ضمنی با معنای تلویحی به معنای عام نیز یکسان نیست. هر استنباط زبانی حامل احساس نیست. تنها زمانی مسئلۀ تحلیل احساس مطرح می‌شود که استنتاج به یک ارزیابی نسبت به هدف منتهی شود. گزارش یک واقعیت می‌تواند صرفاً اطلاع‌رسان باشد، مگر آنکه بافت یا دانش عمومی مطلوبیت آن را مشخص کند.

\subsection{چالش‌های تحلیل احساسات ضمنی}

نخستین چالش، بازیابی شواهدی است که به‌صورت واژۀ احساسی ظاهر نشده‌اند. این شواهد ممکن است در پیامد یک رویداد، برآورده‌شدن انتظار یا رابطۀ علت و معلولی قرار داشته باشند. سامانه باید تشخیص دهد که کدام بخش متن برای هدف اهمیت ارزیابانه دارد.

چالش دوم، استفاده از دانش عمومی و دانش وابسته به دامنه است. یک ویژگی در دو بافت می‌تواند ارزش متفاوتی داشته باشد. برای مثال، «سبک‌بودن» برای یک وسیلۀ قابل حمل ممکن است مطلوب باشد. همین ویژگی برای محصولی که استحکام معیار اصلی آن است، معنای دیگری دارد. در نتیجه، نگاشت ثابت واژه به قطبیت کافی نیست.

چالش سوم، رعایت دامنۀ هدف است. حضور چند موجودیت یا چند جنبه می‌تواند باعث جابه‌جایی قطبیت شود. مدل باید شواهد مربوط به هدف تعیین‌شده را از ارزیابی سایر موجودیت‌ها جدا کند. این مسئله در جمله‌های مقایسه‌ای و جمله‌های دارای چند بند اهمیت بیشتری دارد.

چالش چهارم، تفکیک احساس ضمنی از توصیف خنثی است. استنتاج بیش‌ازحد می‌تواند گزارۀ بی‌طرف را به ارزیابی تبدیل کند. در مقابل، اتکای بیش‌ازحد به نبود واژۀ احساسی باعث ازدست‌رفتن نگرش واقعی می‌شود. تصمیم درست به میزان شواهد و رابطۀ آن با هدف وابسته است.

ابهام، کنایه، وابستگی به فرهنگ و تفاوت انتظار کاربران نیز تحلیل را دشوار می‌کنند. افزون بر این، یک توضیح زبانی روان لزوماً نشان نمی‌دهد که استنتاج درست است. این ویژگی، ارزیابی جداگانۀ مسیر استدلال و برچسب نهایی را ضروری می‌کند.

\section{مدل‌های زبانی بزرگ و مهندسی پرامپت}

مدل زبانی توزیع احتمال دنباله‌های زبانی را بر اساس بافت برآورد می‌کند. مدل‌های زبانی بزرگ\LTRfootnote{Large Language Models (LLMs)} با آموزش روی حجم گسترده‌ای از متن، الگوهای نحوی و معنایی را بازنمایی می‌کنند. بخشی از دانش عمومی نیز در پارامترهای این مدل‌ها بازنمایی می‌شود. مدل زبانی می‌تواند با دریافت یک دستور متنی، پاسخ یا برچسب تولید کند.

\textbf{پرامپت}\LTRfootnote{Prompt} متنی است که وظیفه، ورودی، محدودیت‌ها و شکل پاسخ مورد انتظار را برای مدل مشخص می‌کند. مهندسی پرامپت به طراحی این اجزا برای هدایت رفتار مدل گفته می‌شود. تغییر جزئی در ترتیب دستور، تعریف برچسب‌ها یا قالب خروجی می‌تواند پاسخ را تغییر دهد. ازاین‌رو، پرامپت بخشی از روش آزمایشی است و باید ثبت و ثابت نگه داشته شود.

\subsection{یادگیری درون‌متنی و پرامپت صفرنمونه}

یادگیری درون‌متنی\LTRfootnote{In-Context Learning} به سازگاری رفتار مدل با دستور و نمونه‌های موجود در همان ورودی اشاره دارد. در این روش، وزن‌های مدل تغییر نمی‌کنند. اطلاعات لازم برای انجام وظیفه از طریق بافت ورودی ارائه می‌شود.

در پرامپت چندنمونه‌ای\LTRfootnote{Few-shot Prompting} چند زوج ورودی و خروجی به مدل نشان داده می‌شود. مدل از الگوی این نمونه‌ها برای پاسخ به ورودی جدید استفاده می‌کند. در پرامپت صفرنمونه\LTRfootnote{Zero-shot Prompting} هیچ نمونۀ حل‌شده‌ای ارائه نمی‌شود و مدل فقط تعریف وظیفه و محدودیت‌های پاسخ را دریافت می‌کند.

صفرنمونه‌بودن به معنای نبود دانش پیشین در مدل نیست. این اصطلاح فقط نحوۀ ارائۀ وظیفه در زمان استنتاج را توصیف می‌کند. همچنین، یادگیری درون‌متنی با ریزتنظیم تفاوت دارد. در ریزتنظیم، پارامترهای مدل با داده‌های وظیفۀ جدید به‌روزرسانی می‌شوند، اما در یادگیری درون‌متنی پارامترها ثابت می‌مانند.

\subsection{خروجی آزاد و خروجی ساختاریافته}

در خروجی آزاد، مدل می‌تواند پاسخ را با هر طول و قالبی تولید کند. این حالت برای توضیح مفصل مناسب است، اما پردازش خودکار آن دشوارتر است. ممکن است برچسب در میان توضیح قرار گیرد یا با صورت نوشتاری متفاوتی ظاهر شود.

در خروجی ساختاریافته، نام فیلدها، مقادیر مجاز و ترتیب پاسخ از پیش تعیین می‌شوند. برای مثال، می‌توان از مدل خواست یک برچسب، سطح اطمینان و توضیح کوتاه را در فیلدهای جداگانه برگرداند. ساختار ثابت، اعتبارسنجی و ذخیرۀ خروجی را آسان‌تر می‌کند.

ساختاریافته‌بودن پاسخ، درستی محتوای آن را تضمین نمی‌کند. خروجی باید از دو جنبه بررسی شود: \textbf{اعتبار نحوی}، یعنی انطباق با قالب؛ و \textbf{اعتبار معنایی}، یعنی تعلق مقادیر به دامنه‌های مجاز. یکنواخت‌سازی برچسب‌ها نیز صورت‌های نوشتاری متفاوت را به مجموعۀ ثابت کلاس‌ها نگاشت می‌کند.

\section{استدلال زنجیره‌ای چندمرحله‌ای}

زنجیرۀ تفکر\LTRfootnote{Chain-of-Thought (CoT)} دنباله‌ای از گام‌های زبانی میانی است که ورودی را به پاسخ نهایی مرتبط می‌کند. ایدۀ اصلی این است که مسئلۀ پیچیده به چند تصمیم کوچک‌تر شکسته شود \cite{wei2022cot}. در پیش‌بینی مستقیم، مدل بدون نمایش گام‌های میانی برچسب را تولید می‌کند. در پیش‌بینی زنجیره‌ای، مدل پیش از برچسب، مؤلفه‌های مؤثر بر تصمیم را بیان می‌کند.

برای تحلیل احساسات ضمنی، تجزیۀ مسئله می‌تواند شامل تعیین جنبۀ مرتبط، استنتاج نظر نهفته و تعیین جهت قطبیت باشد. چارچوب \lr{THOR} این منطق را به‌صورت یک استدلال سه‌گامی برای جنبه، نظر و قطبیت صورت‌بندی می‌کند \cite{fei2023thor}. فصل سوم این روش را در بستر پژوهش‌های پیشین بررسی می‌کند و فصل چهارم نحوۀ استفاده از آن را در سامانۀ حاضر توضیح می‌دهد.

مزیت مفهومی استدلال چندمرحله‌ای، ایجاد یک ردپای قابل بررسی است. این ردپا می‌تواند نشان دهد مدل کدام جنبه یا پیامد را مبنای تصمیم قرار داده است. همچنین، خروجی میانی امکان مقایسۀ مسیرهای متفاوت و تشخیص محل احتمالی خطا را فراهم می‌کند.

بااین‌حال، زنجیرۀ طولانی‌تر همواره به پاسخ درست‌تر منتهی نمی‌شود. اگر جنبه اشتباه انتخاب شود، استنتاج‌های بعدی ممکن است بر همان خطا بنا شوند. یک زنجیرۀ منسجم نیز می‌تواند بر فرضی نادرست تکیه کند. بنابراین، روانی توضیح، جانشین ارزیابی برچسب نهایی نیست.

ردپای تولیدشده را نیز نباید گزارش قطعی از سازوکار درونی مدل دانست. این ردپا یک خروجی زبانی قابل مشاهده است که برای تحلیل تصمیم به کار می‌رود. ارزش آن به ارتباط با شواهد متن، سازگاری گام‌ها و انطباق نتیجه با برچسب وابسته است.

\section{بازبینی استدلال و تشخیص خطا}

بازبینی استدلال فرایندی است که پاسخ اولیه را پیش از تصمیم نهایی دوباره بررسی می‌کند. ورودی بازبین می‌تواند متن، هدف، برچسب اولیه و ردپای استدلال باشد. بازبین باید ناسازگاری شواهد، محدودۀ هدف و نتیجه را تشخیص دهد و در صورت لزوم اصلاحی پیشنهاد کند.

بازبینی آزاد معمولاً یک نقد یا پاسخ بازنویسی‌شده تولید می‌کند. بازبینی ساختاریافته، مسئله را به فیلدهایی مانند وضعیت خطا، نوع کلی خطا، برچسب پیشنهادی و سطح اطمینان تقسیم می‌کند. صورت دوم برای تحلیل خودکار مناسب‌تر است، زیرا می‌توان هر فیلد را جداگانه اعتبارسنجی کرد.

تشخیص خطا با اصلاح خطا یکسان نیست. ممکن است بازبین وجود ناسازگاری را درست تشخیص دهد، اما برچسب جایگزین نادرستی پیشنهاد کند. همچنین، نبود خطای اعلام‌شده به معنای صحت پاسخ اولیه نیست. به همین دلیل، خروجی بازبین باید یک منبع اطلاعاتی در فرایند تصمیم‌گیری باشد، نه حقیقت قطعی.

سطح اطمینان نیز لزوماً احتمال کالیبره‌شده نیست. اگر مدل از برچسب‌هایی مانند کم، متوسط و زیاد استفاده کند، این مقادیر فقط ارزیابی خود مدل را نشان می‌دهند. رابطۀ آن‌ها با احتمال تجربی صحت باید جداگانه سنجیده شود.

بازبینی می‌تواند خطای انتقال‌یافته از گام‌های استدلال را آشکار کند، اما خطر بیش‌اصلاحی نیز وجود دارد. در بیش‌اصلاحی، پاسخ درست اولیه فقط به دلیل تولید یک نقد ظاهراً قانع‌کننده تغییر می‌کند. بنابراین، پذیرش اصلاح باید تابع یک سیاست مشخص باشد. ایدۀ افزودن مرحلۀ تأمل به تحلیل احساسات ضمنی در پژوهش‌های مرتبط نیز بررسی شده است \cite{duan2024isa}.

\section{خودسازگاری و رأی اکثریت}

خروجی مدل زبانی در نمونه‌گیری غیرقطعی ممکن است میان اجراها تغییر کند. خودسازگاری روشی است که یک مسیر استدلال را چند بار اجرا و پاسخ نهایی را از میان خروجی‌ها تجمیع می‌کند \cite{wang2023selfconsistency}. این روش، مرحلۀ استدلال تازه‌ای نیست؛ همان روش پیشین را تکرار می‌کند و سپس از رأی خروجی‌ها بهره می‌گیرد.

فرض شود مسیر استدلال $K$ بار اجرا شده و برچسب اجرای $i$ام برابر $y^{(i)}$ است. رأی اکثریت برچسبی را انتخاب می‌کند که بیشترین تعداد تکرار را دارد:

\begin{equation}
    \hat{y}_{\mathrm{vote}}
    = \underset{c \in \mathcal{Y}}{\operatorname{arg\,max}}
    \sum_{i=\text{\lr{1}}}^{K} \mathbb{I}\!\left(y^{(i)}=c\right).
    \label{eq:majority-vote}
\end{equation}

در برابری \ref{eq:majority-vote}، تابع $\mathbb{I}$ در صورت درست‌بودن شرط مقدار یک و در غیر این صورت مقدار صفر دارد. اگر چند کلاس تعداد رأی یکسان داشته باشند، سیاست رفع تساوی باید پیش از ارزیابی تعیین شود. جزئیات این سیاست به پیاده‌سازی وابسته است و در فصل چهارم بیان می‌شود.

نسبت رأی غالب نیز می‌تواند به‌عنوان نشانۀ توافق استفاده شود:

\begin{equation}
    r_{\mathrm{vote}}
    = \frac{\text{\lr{1}}}{K}
    \max_{c \in \mathcal{Y}}
    \sum_{i=\text{\lr{1}}}^{K} \mathbb{I}\!\left(y^{(i)}=c\right).
    \label{eq:vote-agreement}
\end{equation}

مقدار بزرگ‌تر $r_{\mathrm{vote}}$ نشان‌دهندۀ توافق بیشتر اجراها است، اما صحت معنایی پاسخ را تضمین نمی‌کند. چند مسیر می‌توانند خطای نظام‌مند یکسانی داشته باشند. همچنین، اگر اجراها تنوع کافی نداشته باشند، تکرار پاسخ اطلاعات تازه‌ای ایجاد نمی‌کند.

\section{کنترل‌گر و انتخاب چندمنبعی}

سامانۀ چندمنبعی بیش از یک پیش‌بینی برای هر ورودی تولید می‌کند. این منابع می‌توانند شامل پیش‌بینی مستقیم، پیش‌بینی استدلالی و پیشنهاد حاصل از تشخیص خطا باشند. مسئلۀ اصلی آن است که هنگام توافق یا اختلاف منابع، پاسخ نهایی چگونه ساخته شود.

\textbf{کنترل‌گر}\LTRfootnote{Controller} مؤلفه‌ای است که سیگنال‌های میانی را بررسی و یک تصمیم هدایت‌کننده تولید می‌کند. این سیگنال‌ها می‌توانند شامل توافق منابع، اعتبار خروجی، نوع کلی خطا یا سطح اطمینان باشند. کنترل‌گر ممکن است بر اساس قواعد ثابت یا الگوهای آموخته‌شده عمل کند.

\textbf{انتخاب‌گر منبع}\LTRfootnote{Source Selector} مشخص می‌کند کدام منبع، برچسب نهایی را تأمین کند. فرض شود $s_j(x,t)$ پیش‌بینی منبع $j$ و $\mathbf{z}$ بردار سیگنال‌های تصمیم باشد. انتخاب‌گر با سیاست $\pi$ منبع نهایی را تعیین می‌کند:

\begin{equation}
    j^{*}=\pi(\mathbf{z}),
    \qquad
    \hat{y}_{\mathrm{final}}=s_{j^{*}}(x,t).
    \label{eq:source-selection}
\end{equation}

تفکیک کنترل‌گر و انتخاب‌گر اهمیت دارد. خروجی کنترل‌گر می‌تواند یک تصمیم میانی یا پیشنهاد اصلاح باشد، اما پیش‌بینی نهایی لزوماً همان خروجی نیست. این جداسازی منشأ پاسخ نهایی را روشن و ارزیابی هر مؤلفه را ممکن می‌کند.

در انتخاب قاعده‌محور، برای وضعیت‌هایی مانند توافق یا اطمینان زیاد، قواعد از پیش تعریف می‌شوند. در انتخاب داده‌محور، عملکرد منابع روی دادۀ آموزش یا اعتبارسنجی برای ساخت سیاست به کار می‌رود. در هر دو حالت، رفتار پیش‌فرض برای وضعیت‌های ناشناخته باید مشخص باشد.

برچسب طلایی آزمون نباید در ساخت یا تنظیم سیاست انتخاب استفاده شود. استفاده از این اطلاعات نشت داده\LTRfootnote{Data Leakage} ایجاد می‌کند و برآورد عملکرد را خوش‌بینانه می‌سازد. سیاست باید پیش از مشاهدۀ پاسخ‌های طلایی آزمون تثبیت شود. فصل چهارم ویژگی‌ها و قواعد واقعی انتخاب‌گر این پژوهش را تشریح می‌کند.

\section{معیارهای ارزیابی}

ارزیابی طبقه‌بندی، برچسب پیش‌بینی‌شده را با برچسب طلایی مقایسه می‌کند. فرض شود مجموعۀ آزمون شامل $N$ نمونه و $C$ کلاس است. برای هر کلاس $c$، سه مقدار زیر تعریف می‌شوند:

\begin{itemize}
    \item $TP_c$: تعداد نمونه‌های کلاس $c$ که به‌درستی در همان کلاس پیش‌بینی شده‌اند؛
    \item $FP_c$: تعداد نمونه‌های کلاس‌های دیگر که به‌اشتباه در کلاس $c$ قرار گرفته‌اند؛
    \item $FN_c$: تعداد نمونه‌های کلاس $c$ که به‌اشتباه در کلاس دیگری قرار گرفته‌اند.
\end{itemize}

ماتریس درهم‌ریختگی\LTRfootnote{Confusion Matrix} تعداد نمونه‌های هر ترکیب برچسب طلایی و برچسب پیش‌بینی‌شده را نمایش می‌دهد. قطر اصلی این ماتریس پیش‌بینی‌های درست را نشان می‌دهد. خانه‌های خارج از قطر نیز نوع جابه‌جایی میان کلاس‌ها را آشکار می‌کنند.

در متون فارسی، برای دو اصطلاح \lr{Accuracy} و \lr{Precision} گاهی ترجمۀ یکسان «دقت» به‌کار می‌رود. برای جلوگیری از ابهام، در این پایان‌نامه \lr{Accuracy} «دقت کلی» و \lr{Precision} «صحت پیش‌بینی» نامیده می‌شود.

\subsection{دقت کلی}

دقت کلی\LTRfootnote{Accuracy} نسبت تعداد پیش‌بینی‌های درست به کل نمونه‌ها است:

\begin{equation}
    \mathrm{Accuracy}
    = \frac{\text{\lr{1}}}{N}
    \sum_{i=\text{\lr{1}}}^{N} \mathbb{I}\!\left(\hat{y}_i=y_i\right).
    \label{eq:accuracy}
\end{equation}

در برابری \ref{eq:accuracy}، $y_i$ برچسب طلایی و $\hat{y}_i$ برچسب پیش‌بینی‌شدۀ نمونۀ $i$ام است. دقت کلی تفسیر ساده‌ای دارد، اما در داده‌های نامتوازن ممکن است عملکرد ضعیف روی کلاس کم‌تعداد را پنهان کند.

\subsection{صحت پیش‌بینی، بازخوانی و امتیاز اف‌یک}

صحت پیش‌بینی\LTRfootnote{Precision} برای کلاس $c$ نشان می‌دهد چه نسبتی از نمونه‌های پیش‌بینی‌شده در این کلاس واقعاً متعلق به آن هستند:

\begin{equation}
    P_c=\frac{TP_c}{TP_c+FP_c}.
    \label{eq:precision}
\end{equation}

بازخوانی\LTRfootnote{Recall} نشان می‌دهد چه نسبتی از نمونه‌های واقعی کلاس $c$ شناسایی شده‌اند:

\begin{equation}
    R_c=\frac{TP_c}{TP_c+FN_c}.
    \label{eq:recall}
\end{equation}

امتیاز اف‌یک\LTRfootnote{F1-score} میانگین همساز صحت پیش‌بینی و بازخوانی است:

\begin{equation}
    F_{\text{\lr{1}},c}=\frac{\text{\lr{2}}P_cR_c}{P_c+R_c}.
    \label{eq:f1-class}
\end{equation}

میانگین همساز به مقدار کوچک‌تر حساس است. ازاین‌رو، امتیاز اف‌یک زمانی بزرگ می‌شود که صحت پیش‌بینی و بازخوانی هر دو مناسب باشند. اگر مخرج یکی از معیارها صفر باشد، مقدار آن مطابق قرارداد ارزیابی صفر در نظر گرفته می‌شود.

\subsection{امتیاز اف‌یک ماکرو}

امتیاز اف‌یک ماکرو ابتدا اف‌یک هر کلاس را جداگانه محاسبه و سپس میانگین حسابی آن‌ها را گزارش می‌کند:

\begin{equation}
    \mathrm{Macro}\text{-}F_{\text{\lr{1}}}
    =\frac{\text{\lr{1}}}{C}\sum_{c=\text{\lr{1}}}^{C}F_{\text{\lr{1}},c}.
    \label{eq:macro-f1}
\end{equation}

در این معیار، هر کلاس مستقل از تعداد نمونه‌هایش وزن یکسان دارد. بنابراین، عملکرد ضعیف روی کلاس کم‌تعداد اثر آشکارتری بر نتیجه می‌گذارد. گزارش هم‌زمان دقت کلی و اف‌یک ماکرو، تصویری کامل‌تر از کیفیت طبقه‌بندی سه‌کلاسه فراهم می‌کند.

معیارهای کمی فقط صحت برچسب نهایی را می‌سنجند. آن‌ها به تنهایی نشان نمی‌دهند که استدلال از شواهد درست استفاده کرده یا انتخاب‌گر منبع مناسبی را برگزیده است. به همین دلیل، فصل پنجم معیارهای کمی را در کنار تحلیل خطا، بررسی منابع منتخب و نمونه‌های کیفی گزارش می‌کند.

\section{جمع‌بندی}

این فصل واژگان و نمادهای لازم برای ادامۀ پایان‌نامه را تعریف کرد. تحلیل هدف‌محور، قطبیت را حاصل رابطۀ متن، هدف، جنبه و نظر می‌داند. در احساس ضمنی، بخشی از این رابطه آشکار نیست و باید با استفاده از بافت و دانش عمومی استنتاج شود.

مدل زبانی می‌تواند پاسخ مستقیم یا ردپای استدلال چندمرحله‌ای تولید کند. بازبینی، خودسازگاری و کنترل‌گر نیز اطلاعات مکملی برای بررسی و تجمیع پاسخ‌ها فراهم می‌کنند. بااین‌حال، هیچ‌یک به تنهایی تضمین‌کنندۀ صحت نیستند و باید جداگانه ارزیابی شوند. معیارهای دقت کلی و اف‌یک ماکرو نیز مبنای مقایسۀ کمی روش‌ها را فراهم می‌کنند. فصل بعد پژوهش‌های مرتبط را بر اساس همین مفاهیم مرور می‌کند.
